% Appendix M: Semantic Infrastructure, Category Theory, and Knowledge Composition

\section*{M.1 Introduction}

How do multiple cognitive systems interact? How do ideas combine? How do
documents merge? How do repositories evolve? The central claim is:
\emph{Knowledge is not fundamentally stored. Knowledge is composed.}
The mathematical language of composition is category theory.

\begin{definition}{Semantic Category}{sem-cat}
  The \emph{semantic category} $\mathbf{Sem}$ consists of:
  \emph{objects} (concepts, documents, theories, software modules) and
  \emph{morphisms} $f : A \to B$ (translation, summarization, refinement,
  compilation, abstraction). Composition satisfies associativity;
  each object has an identity morphism.
\end{definition}

\begin{definition}{Semantic Merge via Pushout}{sem-merge}
  Given modules $A$ and $B$ sharing interface $I$:
  $A \leftarrow I \to B$, the \emph{merge} is the pushout
  $P = A \sqcup_I B$. This formalizes document merging,
  repository merging, and knowledge integration.
\end{definition}

\begin{definition}{Semantic Conflict}{sem-conflict}
  A \emph{semantic conflict} occurs when no admissible pushout exists:
  contradictory definitions, incompatible assumptions, conflicting ontologies.
  Conflict becomes a mathematical obstruction.
\end{definition}

\begin{definition}{Admissible Merge}{adm-merge}
  A merge is \emph{admissible} if $A(P) \geq \tau$.
  Even if a pushout exists, the result may be incoherent;
  admissibility filters poor merges.
\end{definition}

\begin{definition}{Semantic Functor}{sem-functor}
  A \emph{semantic functor} $F : \mathbf{Sem}_1 \to \mathbf{Sem}_2$
  preserves compositional structure, transporting meaning between domains
  (mathematics~$\to$~software, theory~$\to$~simulation, gesture~$\to$~symbol).
\end{definition}

\begin{definition}{Version History}{version-hist}
  A \emph{version history} is a morphism sequence
  $V_0 \to V_1 \to V_2 \to \cdots$. History becomes categorical;
  traditional version control is a special case.
\end{definition}

\begin{theorem}{Coherent Knowledge Systems are Categories}{knowledge-cat}
  A coherent knowledge system is a category whose objects admit
  admissible composition.
\end{theorem}
\begin{proof}
  Objects alone are insufficient; morphisms provide transformation;
  composition provides growth; admissibility guarantees coherence.
\end{proof}

\begin{namedthm}{The Semantic Infrastructure Principle}
  Knowledge evolves through admissible composition:
  \[
    \text{Objects} \to \text{Morphisms} \to \text{Composition}
    \to \text{Accessibility} \to \text{Civilization.}
  \]
  A civilization is not merely a population. It is a semantic
  infrastructure that preserves, composes, and transmits accessibility
  across generations.
\end{namedthm}
